\documentclass[11pt,letterpaper]{report}
\usepackage[margin=1in]{geometry}
\usepackage[T1]{fontenc}
\usepackage[utf8]{inputenc}
\usepackage{lmodern}
\usepackage{xcolor}
\usepackage{graphicx}
\usepackage{booktabs}
\usepackage{longtable}
\usepackage{array}
\usepackage{enumitem}
\usepackage{fancyhdr}
\usepackage{titlesec}
\usepackage{setspace}
\usepackage[colorlinks=true,linkcolor=black,urlcolor=ttured,citecolor=black]{hyperref}
\usepackage{tcolorbox}
\usepackage{tabularx}

% ---- TTU palette: black, white, minimal red ----
\definecolor{ttured}{RGB}{204,0,0}
\definecolor{ink}{RGB}{17,17,17}
\definecolor{smoke}{RGB}{245,245,245}
\definecolor{mute}{RGB}{100,100,100}

\setlength{\parindent}{0pt}
\setlength{\parskip}{0.6em}
\onehalfspacing

% chapter styling
\titleformat{\chapter}[display]
  {\normalfont\huge\bfseries}{\filleft\large\textcolor{ttured}{PART \thechapter}}{10pt}{\Huge}
\titleformat{\section}{\normalfont\Large\bfseries}{\thesection}{0.8em}{}
\titleformat{\subsection}{\normalfont\large\bfseries}{\thesubsection}{0.7em}{}

\pagestyle{fancy}
\fancyhf{}
\fancyhead[L]{\small\textcolor{mute}{TTU Online Admissions Portal}}
\fancyhead[R]{\small\textcolor{mute}{Technical \& Operational Documentation}}
\fancyfoot[C]{\small\thepage}
\renewcommand{\headrulewidth}{0.4pt}

\newtcolorbox{keybox}[1]{colback=smoke,colframe=ink,title={#1},fonttitle=\bfseries\small,boxrule=0.8pt,arc=1.5mm,left=6pt,right=6pt,top=5pt,bottom=5pt}
\newtcolorbox{warnbox}{colback=red!3,colframe=ttured,title={Important},fonttitle=\bfseries\small,boxrule=0.8pt,arc=1.5mm,left=6pt,right=6pt,top=5pt,bottom=5pt}
\newtcolorbox{plainbox}{colback=white,colframe=mute!40,boxrule=0.6pt,arc=1.5mm,left=6pt,right=6pt,top=5pt,bottom=5pt}

\title{\Huge\bfseries TTU Online Admissions Portal\\[0.4em]
\Large Complete System Documentation\\[0.3em]
\normalsize Architecture, Operations, Data Handling \& Security}
\author{Prepared for April Hollingsworth\\Texas Tech University --- eLearning \& Academic Partnerships}
\date{\today}

\begin{document}

% ---------- TITLE PAGE ----------
\begin{titlepage}
\centering
\vspace*{2cm}
{\color{ttured}\rule{\textwidth}{2pt}}\\[1.2em]
{\Huge\bfseries TTU Online\\[0.15em]Admissions Portal}\\[1em]
{\Large Complete System Documentation}\\[0.5em]
{\color{ttured}\rule{\textwidth}{2pt}}\\[2em]
{\large Architecture, Daily Operations, Data Handling,\\ Identity Matching, and Cybersecurity Measures}\\[2.5em]
\begin{tabular}{ll}
\textbf{Prepared for} & April Hollingsworth \\
\textbf{Department}   & TTU Online, eLearning \& Academic Partnerships \\
\textbf{Platform}     & Web application (mobile-first, desktop-ready) \\
\textbf{Database}     & PostgreSQL (\texttt{lumen\_warehouse}) \\
\textbf{Deployment}   & \texttt{tosmonline0001.ttu.edu} \\
\textbf{Version}      & 1.0 --- \today \\
\end{tabular}
\vfill
{\color{mute}\small Internal document. Contains operational details --- distribute on a need-to-know basis.}
\end{titlepage}

\tableofcontents
\newpage

% ============================================================
\chapter{What This System Is}

\section{The Big Picture}

The TTU Online Admissions Portal is a web application that connects three groups
of people and one very large database:

\begin{itemize}[leftmargin=1.4em]
  \item \textbf{Prospective students} --- people who filled out a ``Request for
        Information'' (RFI) form at an event, on the phone, or online. Their
        details already live in our Salesforce system.
  \item \textbf{Applicants} --- students who create an account in this portal and
        actually apply to a TTU Online program, upload documents, and track
        their progress.
  \item \textbf{Staff} --- advisors who review transcripts and documents,
        managers who combine duplicate records, and administrators who oversee
        the entire operation.
  \item \textbf{The Salesforce mirror and its clone} --- a copy of our
        Salesforce data lives in our own PostgreSQL database
        (\texttt{lumen\_warehouse}) as the \texttt{sf\_origin} schema. A sync
        engine inside the portal re-copies just the TTU~Online slice into the
        app's own \texttt{AH.sf\_*} clone tables every 5--30 minutes. The
        portal reads only its clone, and --- very deliberately --- never writes
        to the mirror automatically.
\end{itemize}

\begin{keybox}{In one sentence}
The portal gives every applicant a clean, mobile-friendly place to apply and
upload documents, while giving staff a back-office control room where every
lead, application, document, and decision is visible --- and where every change
bound for Salesforce is queued for a human to approve by hand.
\end{keybox}

\section{The Problem It Solves}

Before this portal, the workflow had a painful gap. A person might give us their
email and phone number at a recruiting event (becoming a ``lead'' in
Salesforce), then months later apply for a program --- possibly with a different
name spelling, a new phone number, or a different email address. Salesforce
would then hold \emph{two} records that belonged to the \emph{same} human, and
staff had to hunt them down manually.

The portal does three things about this:

\begin{enumerate}[leftmargin=1.6em]
  \item When a student signs up, it \emph{automatically} looks for their
        existing Salesforce record. A perfect email \emph{and} phone match links
        instantly. A partial match goes to a staff review queue.
  \item Staff get a dedicated ``match'' screen to confirm or reject those links,
        and to manually tie a portal account to a Salesforce record that the
        computer could not find.
  \item \textbf{No automatic changes to Salesforce.} Every merge, link, or
        status change lands in a supervised ``outbox'' that a superadmin
        reviews and applies one by one.
\end{enumerate}

\section{What It Looks Like}

The interface is deliberately minimal: black, white, and Texas Tech red used
only as an accent. Sharp corners, strong typography, and big readable tap
targets. It is designed phone-first --- most prospective students will open it
on a phone --- but adapts to a full sidebar layout on desktop monitors.

Two personality touches keep it from feeling like a sterile government form:

\begin{itemize}[leftmargin=1.4em]
  \item \textbf{Raider Red} appears on the student dashboard in three moods:
        sitting patiently when you haven't applied yet, pointing you forward
        while you're working through requirements, and dancing under a rain of
        digital confetti when everything is done.
  \item \textbf{Raider Red is also the chat assistant} --- when a student asks a
        question and no advisor has claimed the conversation yet, a bot persona
        answers and flags staff to jump in.
\end{itemize}

\chapter{The People and Their Roles}

\section{Four Levels of Access}

Every person who logs in belongs to exactly one role. The role decides which
screens exist for them and which actions they can perform.

\begin{center}
\begin{tabularx}{\textwidth}{lX}
\toprule
\textbf{Role} & \textbf{What they can do} \\
\midrule
\textbf{Student} & Apply, upload documents, check requirements, message staff,
                   see their own linked university record, update their profile. \\
\textbf{Advisor} & Everything visible in the admin area that involves reading:
                   leads list, applications list, person profiles, document
                   review (approve / reject / waive), and live chat. \\
\textbf{Manager} & Everything an advisor can do, \emph{plus} the match \& merge
                   screens, status changes (which queue Salesforce work), and
                   analytics. \\
\textbf{Administrator} & Everything: staff account creation, role changes,
                   disabling accounts, the full audit trail, the Salesforce
                   outbox, and user management. \\
\bottomrule
\end{tabularx}
\end{center}

\begin{keybox}{Why roles matter}
An advisor who only reviews transcripts has no business merging identity records
or reading the audit log. Role-based access means each employee sees the
smallest set of tools needed for their job --- which shrinks both mistakes and
attack surface if an account is ever compromised.
\end{keybox}

\section{How Enforcement Works}

Roles are enforced in \emph{two} independent places, not just one:

\begin{enumerate}[leftmargin=1.6em]
  \item \textbf{At the front door.} Every request to an administrative page is
        checked before any page code runs at all. A student's browser never even
        receives the administrative page content --- they are simply redirected
        to their own portal.
  \item \textbf{At every action.} Every request that reads or changes data checks
        the caller's role again. Even if someone crafted a hidden link or typed
        a URL directly, the server answers \emph{``forbidden''} (HTTP 403).
\end{enumerate}

This was tested deliberately: advisors attempting manager-only actions, and
students attempting any administrative action, are all blocked.

\chapter{The Student Journey, Step by Step}

\section{Step 1 --- Signing Up}

A new student visits the portal and creates an account with their email, phone
number, name, and a password. Three things happen in that single moment:

\begin{enumerate}[leftmargin=1.6em]
  \item \textbf{The account is created} in our own project tables (a dedicated
        area of the database called the \texttt{AH} schema --- the project's
        private workspace inside the shared warehouse).
  \item \textbf{The password is transformed} into a one-way cryptographic hash.
        The real password is never stored anywhere --- only a mathematical
        fingerprint that cannot be turned back into a password.
  \item \textbf{The identity match runs.} The system quietly compares the email
        and phone number against nearly fifty thousand sampled Salesforce
        contacts (and millions in production) to answer one question:
        \emph{have we heard from this person before?}
\end{enumerate}

\section{Step 2 --- The Identity Match}

This is the heart of the system. The email a student signs up with may or may
not be the one they gave us as a lead --- so at signup the matcher scores the
new account against the cloned contact slice on four independent signals:
\textbf{email} (40 pts, checked against every email field on the SF record),
\textbf{phone} (35, against every phone field), \textbf{last name} (15,
including former last names --- people change surnames), and \textbf{first
name} (10, including nicknames and preferred names). The total is a
\textbf{confidence score} from 0 to 100. (The SF record's mailing address is
shown to the reviewer as context but isn't scored --- the portal doesn't
collect one yet.) Three outcomes are possible:

\begin{center}
\begin{tabularx}{\textwidth}{llX}
\toprule
\textbf{Match} & \textbf{Result} & \textbf{What the student sees} \\
\midrule
Email + phone + a name all match & \textbf{Auto-linked} (\texttt{auto\_full}) &
Their existing university record appears on the dashboard immediately. \\
Email + phone match, name doesn't & \textbf{Auto-linked + flagged}
(\texttt{auto\_pair}) & Works immediately, but the link sits in a flagged lane
of the match queue until a human eyeballs it --- a name that disagrees is how
duplicates sneak through. \\
Exactly one field hits & \textbf{Pending review} (\texttt{auto\_partial}) &
``We're verifying your record''; a manager picks the right candidate from the
scored list (every candidate, not just the top one). \\
Nothing matches & \textbf{Unlinked} & Normal portal; the match queue offers a
search that scans \emph{all} warehouse contacts, not just the online slice ---
a student's lead may predate it. \\
\bottomrule
\end{tabularx}
\end{center}

\begin{plainbox}
\textbf{Example.} Maria gave us \texttt{maria.g@gmail.com} and phone
(806)\,555-0117 at a college fair last spring. This fall she signs up with the
same email but a new phone number. The system finds the email match, flags the
phone mismatch, and places the link in the review queue. A manager looks at both
records side by side, calls Maria to confirm, and clicks \emph{Confirm}. Only
then is the link made permanent --- and a work item appears in the outbox so a
superadmin can reconcile the duplicate inside Salesforce itself.
\end{plainbox}

\section{Step 3 --- The Application (RaiderConnect)}

The actual application --- demographics, disciplinary questions, program and
term selection --- is completed on \textbf{RaiderConnect}
(\texttt{texastech.my.site.com/online}), Texas Tech's official applicant portal.
That is deliberate: Salesforce owns the application record, and RaiderConnect
is where it must be created.

RaiderConnect cannot be embedded in our page --- Salesforce responds with
\texttt{X-Frame-Options: SAMEORIGIN} and a \texttt{frame-ancestors} CSP that
forbids framing --- so the portal hands the student over in the same tab and
keeps the seam invisible:

\begin{enumerate}[leftmargin=1.6em]
  \item The student picks \textbf{modality} (8-week or 16-week courses ---
        16-week is Education only) and \textbf{program} from the 37-entry
        catalog mirrored from RaiderConnect's own picker, plus a start term.
        This saves an \emph{intent} row (\texttt{rc\_started}) and a timeline
        event.
  \item \emph{Continue to RaiderConnect} takes them to the official form in the
        same tab, with on-page guidance that mirrors the picks they just made
        (``choose Undergraduate $\rightarrow$ Fully Online $\rightarrow$
        8-week $\rightarrow$ \emph{their program}'').
  \item When they come back --- detected automatically on tab focus, a 30-second
        poll, or the \emph{I've submitted --- check now} button --- the portal
        looks for a new TTU~Online application on their linked contact, live in
        the warehouse first (fresher than the clone). When it lands, the intent
        row attaches the Salesforce application id and flips to
        \texttt{submitted}.
\end{enumerate}

\section{Step 4 --- The Checklist}

Once the Salesforce application exists, the checklist \emph{is Salesforce's}:
every \texttt{TargetX\_SRMb\_\_Checklist\_Item\_\_c} row for that application is
piped into the warehouse (\texttt{AH.checklist\_items}, via the Lumen ``SF
Checklist Items'' connector, incremental) and cloned into
\texttt{AH.sf\_checklist\_item} each sync --- name, requirement code, status,
received/closed/waived flags. What a student uploaded on RaiderConnect shows up
received on the next sync, automatically.

The local \texttt{portal\_checklist\_items} template checklist remains only as a
fallback while no SF application exists yet --- and for \textbf{General
Business} applicants the portal shows ``requirements coming soon,'' because
that program's requirements are still being finalized on the TTU side.

A student can also upload straight to the portal for an open requirement; that
upload counts as \emph{submitted with us} until Salesforce parity confirms it
(next section).

\section{Why the Checklist Is Built This Way --- the psychology}

Because TTU Online charges no application fee, nothing financial stops an
applicant from drifting away mid-checklist --- so the checklist is designed
around four well-studied behavioral effects rather than around reminders that
get ignored.

\begin{itemize}[leftmargin=1.4em]
  \item \textbf{Endowed progress} (Nunes \& Dr\`eze, \emph{Journal of Consumer
        Research}, 2006 --- the car-wash study: a ten-stamp card with two free
        stamps fills faster than an eight-stamp card). Progress that
        \emph{appears already started} reframes the task as ``in progress,''
        not ``not begun.'' A new student's checklist therefore never starts at
        zero: ``Create your portal account'' and ``Submit your application''
        are real milestones counted in the same denominator as the requirement
        rows, so a signed-up applicant opens the page at \texttt{2 of 6}, not
        \texttt{0 of 4}. Both milestones reflect genuine completed work --- no
        fabricated credit.
  \item \textbf{Goal gradient} (Kivetz, Urminsky \& Zheng, \emph{Journal of
        Marketing Research}, 2006 --- people accelerate as a goal gets closer).
        The header copy changes with distance: ``$n$ of $m$ complete'' early,
        ``1 item left --- you're almost there'' at the end, and the single next
        action is surfaced beside the progress bar so the closest item gets the
        attention.
  \item \textbf{Open loops} (the Zeigarnik effect, 1927 --- unfinished tasks
        stay mentally active, but only while they're visible). Remaining items
        stay on screen; items confirmed received by Texas Tech trigger a
        one-time ``Texas Tech received\ldots'' callout so completion feels like
        an event. Guardrail: this effect is debated in the literature and only
        honest remaining work is ever shown --- fake ``80\% done'' loops are a
        dark pattern and were deliberately not used.
  \item \textbf{Personalized, actionable nudges} --- Castleman \& Page's
        ``summer melt'' texting studies lifted enrollment with status-aware
        messages carrying a concrete next step, while large randomized trials
        (Oreopoulos et al.; the College Board RYCP trial on $\sim$785{,}000
        students) found generic informational blasts produce roughly nothing.
        The practical rule the portal follows: every nudge names the specific
        outstanding item and a direct action, and stalled students surface to a
        human in Tickets rather than receiving automation-only contact.
        Literal daily streaks and loss-aversion mechanics (Duolingo-style)
        were evaluated and rejected --- admissions has no daily action to
        streak on, and punishing a stalled applicant only increases
        abandonment.
\end{itemize}

\subsection{What this looks like in the interface}

A slim progress bar sits atop the checklist page with the honest count
(\texttt{2 of 6}), the next action named beside it, and milestone callouts
when the sync engine flips a requirement to received. No gamification chrome
beyond that --- the design stays minimal on purpose, since clutter itself is a
documented abandonment driver.

\subsection{Deletion and data minimization}

Profile deletion is real deletion, not deactivation: the \texttt{portal\_users}
row is removed, cascading through identity links, applications, checklist
items, documents (files are unlinked from disk first), chat threads, and the
contact projection; activity-log rows keep their anonymized shell via
\texttt{ON DELETE SET NULL} so the audit trail survives without retaining the
person. The student's official Texas Tech record in Salesforce is untouched ---
the portal was never the system of record. Consistent with the terms page, no
data beyond first name, email, phone, and uploaded documents is collected, and
no feature asks for more.

\section{Step 5 --- Documents}

Students upload files (PDF, JPG, PNG, HEIC, DOC, DOCX --- up to 15\,MB each),
optionally attached to a specific checklist item. Every upload is:

\begin{itemize}[leftmargin=1.4em]
  \item \textbf{Verified by content, not by name} --- the server inspects the
        file's actual bytes (its ``magic bytes'') rather than trusting what the
        phone claims the file type is.
  \item \textbf{Renamed to a random identifier} on disk, so an uploaded file can
        never be executed as a script or reached by guessing a URL.
  \item \textbf{Encrypted before it hits the disk} (AES-256-GCM --- the same
        class of encryption used by banks) and decrypted only in memory when an
        authorized person opens it.
  \item \textbf{Access-checked} --- a student can open only their own files;
        staff roles can open anyone's.
\end{itemize}

\emph{Document status is a parity loop, not a single flag:}

\begin{center}
\texttt{pending\_review} $\rightarrow$ \texttt{pushed} $\rightarrow$
\texttt{approved} \qquad (\texttt{rejected} is the sideways exit)
\end{center}

\begin{itemize}[leftmargin=1.4em]
  \item \textbf{pending\_review} --- lands here the moment it's uploaded; appears
        in the staff document queue (\texttt{/admin/documents}).
  \item \textbf{pushed} --- an advisor downloaded the file and attached it to
        the student's checklist inside RaiderConnect/Salesforce by hand, then
        marked it pushed (optionally picking which checklist item it satisfies).
        An outbox ticket records who did it.
  \item \textbf{approved} --- set \emph{only} by the sync engine: once the SF
        checklist item shows received/closed/waived, parity is real and the
        portal calls it approved. Nobody marks this by hand.
\end{itemize}

\section{Step 6 --- Messaging}

A persistent chat thread connects the student to staff. When a student writes
and nobody has claimed the conversation, Raider Red (the bot persona) answers
immediately with a helpful holding message --- or, once the custom knowledge-base
bot is connected, a real answer. The moment an advisor replies, the thread is
``claimed'' and the human takes over. Staff see a shared queue of open threads
on their side.

\section{Step 7 --- Sharing the Moment}

Once an application exists, a \emph{``Tell everyone''} button appears. On a phone
it opens the native share sheet (so the student can post ``I just applied to TTU
Online!'' to Instagram or text it to family); on a computer it copies a link to
the clipboard.

\chapter{The Staff Side: The Admin Portal}

\section{Dashboard}

The admin landing page is a control-room summary: counts of students,
applications, leads, pending reviews, unread chat, and a live feed of the most
recent events across the whole system.

\section{Leads}

The leads list is a searchable, filterable window onto every RFI contact in the
cloned slice --- the same records the recruiting team creates in Salesforce,
at most one sync interval old. Staff search by name, email, or phone, and
filter by lifecycle stage (Prospect, Inquiry, Applicant\ldots) and term.

\section{Applications}

Every application known to the system --- both the ones students filed through
this portal and the ones already sitting in Salesforce --- with status, stage,
term, program, and record type. Clicking any row opens the person view.

\section{Person 360}

One page per human. It combines:

\begin{itemize}[leftmargin=1.4em]
  \item Salesforce contact fields (name, email, phone, status, term),
  \item every application on file and its stage,
  \item every recorded interaction (RFI events, touchpoints),
  \item the linked portal account, if one exists,
  \item the portal application's \textbf{status timeline} --- every transition
        with actor and timestamp, from draft through decision,
  \item uploaded documents and checklist state,
  \item staff notes,
  \item and that person's own activity on the website --- which pages they
        visited and how much time they spent.
\end{itemize}

This is where an advisor can answer, in seconds, ``who is this person and where
are they in the process?''

\section{Match \& Merge}

The manager-facing screen has three lanes:

\begin{enumerate}[leftmargin=1.6em]
  \item \textbf{Pending review} --- partial matches awaiting a human yes/no,
        with both records shown side by side and exactly which fields matched
        highlighted.
  \item \textbf{Unlinked users} --- portal accounts with no Salesforce match,
        each with the system's best-guess candidate contacts attached.
  \item \textbf{Auto-linked} --- the perfect matches that linked themselves, so
        staff can spot-check that automation behaved.
\end{enumerate}

Every suggested pair carries a \textbf{confidence score} (0--100) and
per-field evidence chips --- email, phone, last name (\emph{including former
last names, which catches post-marriage surname changes}), and first name ---
while the contact's city/state sits alongside as extra context, so the
reviewer sees \emph{why} the system suspects a match, not just that it does.
A free-text search over the clone covers the case where the algorithm missed
entirely.

\begin{warnbox}
Confirming a link or requesting a merge does \textbf{not} modify Salesforce.
It writes a work item into the outbox. A superadmin applies each change by
hand, one at a time, inside Salesforce proper.
\end{warnbox}

\section{Review Queue}

A single inbox of everything waiting on a human: submitted checklist items and
freshly uploaded documents. Approve, reject (with a note the student sees), or
waive --- every click is recorded with the advisor's name.

\section{Status Changes and the Outbox}

When a manager changes an application's status --- say, from \emph{submitted} to
\emph{under review} or \emph{admitted} --- three writes happen:

\begin{enumerate}[leftmargin=1.6em]
  \item the portal's own copy is updated so the student sees progress instantly;
  \item a row is appended to \texttt{portal\_application\_events} --- the
        append-only timeline recording \emph{from-status, to-status, actor,
        timestamp}, so an application's full progression is never lost;
  \item a row is appended to the \textbf{outbox} describing exactly what must be
        mirrored into Salesforce.
\end{enumerate}

The outbox is a screen of pending work items for the superadmin: entity, action,
payload, who requested it, and when. Each item can be marked \emph{done} or
\emph{skipped} with a note. This is the deliberate air-gap between the public
website and the institutional system of record: \textbf{the website cannot
change Salesforce; only a human can.}

\section{Live Chat}

The staff console lists every open student conversation, shows the last message
and whether Raider Red is still holding the thread, and lets any advisor claim,
reply, or close. Replying automatically claims the thread.

\section{Users \& Roles}

Administrators create staff accounts, change roles, and disable accounts. A
disabled account is locked out immediately at every layer.

\section{Audit Trail}

Every meaningful thing that happens is written to an append-only activity table
with the actor, the action, the detail, the IP address, the browser, and the
timestamp:

\begin{itemize}[leftmargin=1.4em]
  \item every page view and click by \emph{both} students and employees,
  \item every login, logout, failed login, MFA failure, and password reset
        (flagged as security events),
  \item every admin action --- reviews, merges, status changes, chat replies,
  \item and time-on-site per user, powering the analytics screen.
\end{itemize}

The audit view is admin-only and exists for exactly the questions ``who did
this?'' and ``what was this employee doing at 2\,pm yesterday?''

\section{Analytics}

A daily event chart, the most-visited pages, and per-user time-on-site --- the
``how much time did a student spend on the site'' requirement made visible.

\chapter{Where the Data Lives}

\section{Three Zones, Three Responsibilities}

Everything sits in one PostgreSQL database called \texttt{lumen\_warehouse},
divided into clearly separated areas (``schemas''):

\begin{center}
\begin{tabularx}{\textwidth}{lX}
\toprule
\textbf{Zone} & \textbf{Contents} \\
\midrule
\texttt{sf\_origin} & The Salesforce mirror: contacts (3.5M), applications
(470K), interactions (4.3M), accounts, terms, programs, record types. Fed by
the warehouse's own ETL. The portal \textbf{reads} it during syncs and
\textbf{never writes} to it. \\
\texttt{AH.sf\_*} & The project's \textbf{clone} of the TTU~Online slice ---
contact, application, term, program, record type --- refreshed every 5--30
minutes, plus append-only \texttt{*\_hist} tables keeping every observed
version of every row. All portal reads hit this copy. \\
\texttt{AH.portal\_*} & The app's own data: users, sessions, identity links,
applications, checklists, documents, chat, notes, audit log, outbox,
projections. The only zone the portal writes interactively. \\
\bottomrule
\end{tabularx}
\end{center}

This separation is the project's most important architectural rule. The portal
can never corrupt institutional data --- its writes are physically confined to
its own tables, and in production the database account is granted
\emph{read-only} permission on \texttt{sf\_origin} as a second,
belt-and-suspenders layer. Keeping our own clone (rather than querying the
mirror live) means the admissions site keeps working even if the warehouse ETL
is down, lets us attach portal-specific columns to records, and gives us a
searchable append-only history of how each record evolved.

\section{The Key Tables, in Plain Language}

\begin{center}
\begin{tabularx}{\textwidth}{lX}
\toprule
\textbf{Table (in \texttt{AH})} & \textbf{What it holds} \\
\midrule
\texttt{portal\_users} & Accounts: email, hashed password, role, MFA flag,
failed-login counter, lock timestamp. \\
\texttt{portal\_sessions} & Active sign-ins: who, when, from which IP/device,
expiry. \\
\texttt{portal\_otp} & Text-message codes: stored as hashes, 10-minute life,
5-try limit. \\
\texttt{portal\_identity\_links} & The bridge: which portal account belongs to
which Salesforce contact, how it matched, who reviewed it. \\
\texttt{portal\_applications} & Applications filed in the portal (program, term,
status, answers). \\
\texttt{portal\_checklist\_items} & Each applicant's requirement list and its
state. \\
\texttt{portal\_checklist\_templates} & The editable requirement lists per
program/term --- where the real admissions requirements will live. \\
\texttt{portal\_documents} & Document metadata (the files themselves live,
encrypted, on disk). \\
\texttt{portal\_chat\_threads/messages} & Conversations with advisors and the bot. \\
\texttt{portal\_outbox} & The queue of changes a superadmin must apply to
Salesforce by hand. \\
\texttt{portal\_activity} & The audit trail: every click, page, login, and admin
action. \\
\texttt{portal\_notes} & Staff notes attached to people. \\
\texttt{portal\_application\_events} & Append-only status timeline per
application --- every transition with actor and timestamp (stateful, not just
the latest status). \\
\texttt{portal\_contact\_projection} & The ``generated columns'': for each
linked SF contact, the portal user id, latest portal application + status,
the full status-history JSON, checklist progress, and last portal activity.
Rebuilt each sync and on every link/status change. \\
\texttt{sf\_contact / sf\_application} & The cloned slice: current state of
the records the portal cares about (denormalized, resolved names). \\
\texttt{sf\_contact\_hist / sf\_application\_hist} & Append-only history: every
version of every cloned row, keyed by content hash --- how an SF record
progressed through statuses over time. \\
\texttt{sf\_term / sf\_program / sf\_recordtype} & Cloned reference data for
the pickers and filters. \\
\texttt{sf\_sync\_runs} & Every sync logged: trigger (schedule/manual/boot),
row counts, duration, errors --- the audit trail for the pipeline itself. \\
\bottomrule
\end{tabularx}
\end{center}

\section{How Sync Works --- Clone Down, Queue Up}

The user asked for one rule above all: \emph{do not let the app mutate the
warehouse on its own}. Sync therefore runs in two opposite, asymmetric
directions:

\textbf{Down (Salesforce $\rightarrow$ portal): automatic, on a timer.}
\begin{enumerate}[leftmargin=1.6em]
  \item A sync engine inside the app wakes every \textbf{5--30 minutes}
        (\texttt{SYNC\_INTERVAL\_MINUTES}, default 15) and once at boot if the
        clone is empty or stale.
  \item Each run takes a Postgres advisory lock first, so a scheduled run and a
        manual \emph{Refresh now} can never overlap.
  \item It re-reads just the TTU~Online slice of the 3.6-million-row mirror ---
        contacts flagged Distance Learning, \texttt{TTU\_Online\_Student},
        online lead sources, TUO campus, or owning a TTU~Online application
        ($\approx$76{,}500 contacts) --- plus their applications ($\approx$33K),
        per-item checklist rows ($\approx$55K, TTU~Online applications only),
        and the small reference tables --- and \textbf{upserts by Salesforce Id}
        into the \texttt{AH.sf\_*} clone.
  \item Every row carries a content hash; when a row differs from last time it
        is \emph{also} appended to \texttt{sf\_*\_hist} --- so the clone is
        current-state \textbf{and} the history is append-only (no state is
        ever lost, on either side).
  \item After the copy, two reconciliation steps run: \emph{detect} attaches a
        newly-found Salesforce application to the student's intent row (and
        writes a \texttt{submitted} timeline event), and \emph{parity} flips
        pushed documents to \texttt{approved} once their SF checklist item
        shows received. Then \texttt{portal\_contact\_projection} is rebuilt:
        the generated columns stitching each linked contact to its portal user,
        application status + timeline, SF-side status/stage/decision, checklist
        counts, open items, and pending uploads.
  \item Every run lands in \texttt{sf\_sync\_runs}. Managers and admins see
        run history and a \textbf{Refresh now} button at \texttt{/admin/sync}.
\end{enumerate}

\textbf{Up (portal $\rightarrow$ Salesforce): manual, always.}
\begin{enumerate}[leftmargin=1.6em]
  \item \textbf{Writes stay local.} Anything that changes state lands in
        \texttt{AH.portal\_*} tables instantly (the portal is always truthful
        about the portal) --- including a \texttt{portal\_application\_events}
        row for every status transition.
  \item \textbf{Salesforce changes queue as tickets.} Status flips, merges, and
        link confirmations each write an outbox row: \emph{entity, action,
        payload, requester, timestamp}. The superadmin works the queue one item
        at a time and marks each done; the next scheduled sync then reflects
        the change back into the clone --- closing the loop without the
        website ever touching Salesforce.
\end{enumerate}

\begin{keybox}{Why a queue instead of a direct sync}
Because identity is messy. If the app merged two Salesforce records
automatically and got it wrong, un-merging inside Salesforce is a nightmare of
re-parenting related objects. A human approving each merge costs minutes and
prevents disasters --- the right trade at 200--300 daily users.
\end{keybox}

\chapter{Cybersecurity, Explained for Non-Engineers}

\section{The Layer Cake}

Security here is not one lock but a stack of independent defenses --- if any one
layer fails, the next still holds. What follows is every layer, in plain terms.

\section{Layer 1 --- Passwords}

\begin{itemize}[leftmargin=1.4em]
  \item Passwords are hashed with \textbf{scrypt} at deliberately expensive
        settings (roughly 128\,MB of memory work per attempt). Guessing a
        password by brute force takes a real fraction of a second \emph{per
        guess} --- fast enough for one login, impossibly slow for billions of
        guesses.
  \item A \textbf{server-side pepper} (a secret string mixed in before hashing,
        held only in server configuration --- never in the database) means that
        even an attacker who steals the entire user table still cannot test
        guesses without also stealing the server environment.
  \item Password rules: at least 10 characters, must include a letter and a
        number, cannot be a common password, cannot contain the user's own
        email.
  \item Old password hashes upgrade silently --- when a user with an older hash
        format logs in successfully, their password is re-hashed at the new
        standard automatically.
  \item Login is \textbf{timing-safe}: the same amount of hashing work happens
        whether or not the email exists, so an attacker cannot probe ``is this
        email registered?'' by measuring response time.
\end{itemize}

\section{Layer 2 --- Multi-Factor Authentication}

\begin{itemize}[leftmargin=1.4em]
  \item MFA uses \textbf{Twilio Verify} --- codes are generated and checked by
        Twilio, the industry-standard service, not by us.
  \item Codes expire in \textbf{10 minutes} and allow at most \textbf{5 attempts}
        per challenge.
  \item Only a \emph{hash} of each code is ever stored in our database.
  \item The portal records every failed MFA attempt as a security event in the
        audit log.
  \item In development (no Twilio keys configured) a fixed code
        \texttt{424242} is used --- production prints a loud startup warning if
        Twilio keys are ever missing.
\end{itemize}

\section{Layer 3 --- Sessions (Staying Logged In)}

\begin{itemize}[leftmargin=1.4em]
  \item A login creates a signed token (like a tamper-proof wristband) stored in
        a cookie that JavaScript cannot read (\texttt{httpOnly}), that only
        travels on encrypted connections in production (\texttt{secure}), and
        that is never sent to other websites (\texttt{sameSite}).
  \item Every session also has a row in the database with its expiry, IP, and
        device --- so sessions can be revoked instantly.
  \item Changing a password \textbf{kills every existing session} for that
        account. A stolen laptop stops being logged in the moment the password
        is reset.
  \item Sessions expire in 7 days regardless.
\end{itemize}

\section{Layer 4 --- Rate Limiting (Brute-Force Armor)}

Every sensitive doorway has a speed limit. Exceed it and the server answers
``slow down'' (HTTP 429) with a retry timer:

\begin{center}
\begin{tabular}{ll}
\toprule
\textbf{Action} & \textbf{Limit} \\
\midrule
Login attempts            & 10 per 5 minutes (per IP), \emph{plus} 5 failures $\Rightarrow$ 15-min account lock \\
Account creation          & 5 per hour (per IP) \\
Password reset requests   & 5 per hour (per IP) + 3 per hour per email \\
MFA code checks           & 10 per 5 minutes \\
Document uploads          & 20 per 10 minutes \\
Chat messages             & 60 per minute \\
Activity tracking pings   & 120 per minute \\
Everything else           & 240 per minute \\
\bottomrule
\end{tabular}
\end{center}

\begin{plainbox}
\textbf{What this prevents.} An attacker who wants to try ten thousand passwords
against one account hits the lock at \emph{five}. Someone trying to create
thousands of fake accounts stops at five per hour per IP. It converts
unlimited-force attacks into mere annoyances.
\end{plainbox}

\section{Layer 5 --- Forgery Protection (CSRF \& CORS)}

\begin{itemize}[leftmargin=1.4em]
  \item Every request that \emph{changes} something must arrive with an
        \texttt{Origin} header matching our own site. A malicious page on
        another domain cannot silently make your browser submit forms here ---
        its requests are refused outright. (Verified: requests forged from a
        foreign origin get a 403.)
  \item No ``CORS'' permissions are granted at all: other websites' JavaScript
        simply cannot read our API responses.
  \item Cross-origin isolation headers (\texttt{Cross-Origin-Opener-Policy},
        \texttt{Cross-Origin-Resource-Policy}) prevent other tabs from sharing a
        window or loading our authenticated resources.
\end{itemize}

\section{Layer 6 --- Security Headers}

Every response carries a set of instructions that harden the browser itself:

\begin{center}
\begin{tabularx}{\textwidth}{lX}
\toprule
\textbf{Header} & \textbf{In plain terms} \\
\midrule
Content-Security-Policy & ``Only run scripts/styles/images from this site ---
nothing external is allowed to execute.'' Blunts injected-script attacks. \\
Strict-Transport-Security & ``Always use HTTPS with this site for the next two
years --- never fall back to plain HTTP.'' \\
X-Frame-Options / frame-ancestors & ``Never let this site be embedded in a
frame'' --- defeats clickjacking overlays. \\
X-Content-Type-Options & ``Don't guess file types'' --- a file we call an image
can never be reinterpreted as code. \\
Referrer-Policy & Pages don't leak full URLs to third parties. \\
Permissions-Policy & Camera, microphone, and location are disabled --- the app
never asks for them. \\
Cache-Control: no-store & Authenticated pages and API answers are never saved
in browser or proxy caches. \\
X-Robots-Tag: noindex & Search engines are told not to index portal/admin/API
pages. \\
\bottomrule
\end{tabularx}
\end{center}

The framework also hides its own name (\texttt{poweredByHeader} is off) ---
there's no reason to advertise the exact software stack to scanners.

\section{Layer 7 --- Injection Defense}

\begin{itemize}[leftmargin=1.4em]
  \item \textbf{SQL injection is structurally impossible} in the current code:
        every query sends user data to the database \emph{separately} from the
        command, as parameters --- the database never interprets a name or
        email as a command.
  \item \textbf{Cross-site scripting (XSS)} is handled by the framework's
        escaping plus the strict CSP above; user text is rendered as text,
        never as HTML.
  \item \textbf{Every input is capped}: emails, names, notes, and chat bodies
        are trimmed to sane lengths so nobody can flood a field with a
        megabyte of junk.
  \item \textbf{Path traversal is double-checked}: document downloads validate
        that the requested ID is a well-formed UUID, that the stored filename is
        one we generated (not something containing ``\texttt{../}''), and that
        the resolved file path really is inside the upload directory.
\end{itemize}

\section{Layer 8 --- Files (Documents)}

The riskiest feature of any admissions portal is file upload. The defenses:

\begin{enumerate}[leftmargin=1.6em]
  \item \textbf{Extension allowlist} --- only PDF, JPG, PNG, HEIC, DOC, DOCX.
  \item \textbf{Magic-byte inspection} --- the first bytes of the file must
        actually be a real PDF/image/document signature. A script renamed
        \texttt{transcript.pdf} is rejected.
  \item \textbf{Random storage names} --- the file lands on disk under a UUID
        filename in a folder outside the public web directory. It can never be
        addressed by URL or executed.
  \item \textbf{Encryption at rest} --- AES-256-GCM before writing; the server
        decrypts only in memory when an authorized viewer opens it.
  \item \textbf{Ownership checks on every download} --- a student can open only
        their own files; staff roles can open all. Verified: a second student
        attempting to open someone else's document gets 403.
  \item \textbf{Size cap} --- 15\,MB, enforced server-side.
  \item \textbf{Rate limit} --- 20 uploads per 10 minutes.
\end{enumerate}

\section{Layer 9 --- Encryption in Transit}

\begin{itemize}[leftmargin=1.4em]
  \item Between browser and server: HTTPS only. The session cookie refuses to
        travel unencrypted, and HSTS pins the site to TLS for two years.
  \item Between application and database: TLS is supported via standard
        PostgreSQL SSL modes (\texttt{sslmode=verify-full}) when the two live on
        separate hosts; on one host, connections never leave the machine.
\end{itemize}

\section{Layer 10 --- Watching Everything (Audit)}

Every login, failed login, MFA failure, password reset, page view, click,
upload, review decision, merge, and status change lands in the audit table with
who/what/where/when. Security-sensitive events are tagged distinctly so the
audit screen can filter for them. This is both a compliance control and a
detection tool --- unusual patterns (many failed logins, an employee suddenly
bulk-viewing records) are visible in one place.

\section{Layer 11 --- Boot-Time Safety Checks}

At startup in production, the server \emph{refuses to boot} if the session
secret is missing or weak, and prints loud warnings if encryption keys, the
password pepper, or Twilio credentials are absent --- so the app can never
silently drift into an insecure configuration.

\section{The Deliberate Air-Gap}

The deepest security property isn't cryptographic at all: \textbf{the app has no
code path that writes to the Salesforce mirror}. It's not that writes are
restricted --- they don't exist. Even a fully compromised application process
cannot alter institutional records; at worst it could queue outbox tickets a
human would review.

\chapter{Running the System Day to Day}

\section{What Staff Do Each Morning}

\begin{enumerate}[leftmargin=1.6em]
  \item \textbf{Advisors} open the document queue (\texttt{/admin/documents}),
        download each \texttt{pending\_review} file, attach it to the right
        checklist item inside RaiderConnect, and mark it \texttt{pushed}
        (optionally naming which SF item it satisfies) --- and watch the chat
        console for new conversations.
  \item \textbf{Managers} check the match queue: confirm or reject pending
        identity links, manually link the stragglers, and move applications
        through their statuses --- knowing every flip also writes an outbox
        ticket and lands on the student's status timeline.
  \item \textbf{Administrators} review the audit trail for anomalies, manage
        staff accounts, glance at \texttt{/admin/sync} to confirm the clone is
        fresh (or force a refresh), and work the outbox: applying each queued
        change in Salesforce by hand and marking it done.
\end{enumerate}

\section{The Normal Data Flow}

\begin{plainbox}
\textbf{End-to-end example --- RFI to admit:}
\begin{enumerate}[leftmargin=1.4em,itemsep=2pt]
  \item Recruiter collects an RFI at an event $\rightarrow$ Salesforce lead
        appears in \texttt{sf\_origin} via the warehouse ETL.
  \item Next sync (or a manager's \emph{Refresh now}) clones the lead into
        \texttt{AH.sf\_contact}; any change since last run is appended to
        \texttt{sf\_contact\_hist}.
  \item Student signs up on their phone $\rightarrow$ the matcher scores them
        against the clone (email/phone/last/first name) $\rightarrow$ auto-link,
        or queued for review with a confidence score.
  \item Student picks program + term and continues to RaiderConnect
        $\rightarrow$ intent row saved (\texttt{rc\_started}); the official
        application is created there, in Salesforce.
  \item The next sync (or the return check) finds the new TTU~Online
        application on their contact $\rightarrow$ the intent row attaches it,
        flips to \texttt{submitted}, timeline event recorded.
  \item Student uploads a transcript in our portal $\rightarrow$ file encrypted,
        status \texttt{pending\_review}, advisor sees it in the document queue.
  \item Advisor attaches it to the checklist item inside RaiderConnect and marks
        the portal copy \texttt{pushed} $\rightarrow$ outbox ticket queued.
  \item Next sync sees the SF checklist item received $\rightarrow$ the portal
        document flips to \texttt{approved} by itself. The record shows
        identical state on both sides, and a human certified every write to
        Salesforce.
\end{enumerate}
\end{plainbox}

\section{Capacity}

The workload assumptions baked into the design (200--300 users/day, 2,000--5,000
per month) are modest by web standards. The database connection pool (8
connections), in-memory rate limiter, and polling chat are sized for exactly
this scale --- nothing in the design bottlenecks until orders of magnitude more
traffic. The one component that \emph{will} want upgrading first at scale is the
chat's 2.5-second polling, which can become server-sent events or websockets
later without changing the user experience.

\chapter{Production --- \texttt{tosmonline0001.ttu.edu}}

\section{The Shape of Production (as deployed)}

\begin{itemize}[leftmargin=1.4em]
  \item \textbf{Two boxes, one public face.} The portal app runs on
        \texttt{tosmonline0001} --- the public-facing host, where nginx already
        terminates :443 with the real TTU certificate. The databases stay on
        \texttt{tosmonline0003}: \texttt{lumen\_warehouse} (SF mirror +
        \texttt{AH.checklist\_items} via the \emph{SF Checklist Items} Lumen
        connector) and \texttt{AHPortalDB} (every \texttt{portal\_*} table plus
        the \texttt{sf\_*} clone, schema \texttt{"AH"}).
  \item \textbf{Cross-box database access} rides a persistent SSH tunnel ---
        Postgres on 0003 binds localhost only, so PM2 process
        \texttt{ahportal-db-tunnel} keeps
        \texttt{127.0.0.1:55432 $\rightarrow$ 0003:5432} open (key auth between
        the boxes already existed). The app's \texttt{DATABASE\_URL} and
        \texttt{WAREHOUSE\_URL} both point at \texttt{127.0.0.1:55432}.
  \item \textbf{Two PM2 processes} on 0001 (\texttt{pm2 save}d): the tunnel and
        \texttt{april-web} (\texttt{next start -p 3210},
        \texttt{NEXT\_BASE\_PATH=/sim/AHPortal}).
  \item \textbf{Public URL} \texttt{https://tosmonline0001.ttu.edu/sim/AHPortal/}:
        an nginx \texttt{location} block in \texttt{chatbot.conf} proxies to
        \texttt{127.0.0.1:3210} and forwards \texttt{X-Forwarded-Host} /
        \texttt{-Proto} / \texttt{-For}, which the app's CSRF check and logs
        consume.
  \item \textbf{Least privilege:} the portal's \texttt{ahportal} role owns
        \texttt{AHPortalDB.AH} and holds \emph{SELECT-only} grants on
        \texttt{lumen\_warehouse.sf\_origin} and \texttt{AH}. Writes to the
        mirror fail at the permission level.
  \item \textbf{Fresh server-side secrets} live in \texttt{.env.local} on 0001;
        uploads sit in \texttt{\textasciitilde/ahportal/uploads}, outside any
        web path.
\end{itemize}

\begin{warnbox}
The box's glibc (2.28) predates what Next.js's native SWC binary needs ---
production builds must run \texttt{next build --webpack} (the npm
\texttt{build} script does this). Public assets must be referenced through
\texttt{asset()} (adds the \texttt{/sim/AHPortal} prefix); bare root-relative
\texttt{/logo.png}-style paths 404 behind the tenant mount.
\end{warnbox}

\begin{warnbox}
\textbf{Before launch:} the dev values currently in \texttt{.env.local}
(passwords, pepper, encryption key) must be regenerated on the server, and the
dev MFA code \texttt{424242} disappears automatically once Twilio keys are set.
\end{warnbox}

\section{Go-Live Checklist (what was actually done)}

\begin{enumerate}[leftmargin=1.6em]
  \item \textbf{Provision:} created \texttt{AHPortalDB} + least-privilege
        \texttt{ahportal} role --- owns the portal schema, \texttt{SELECT}-only
        on \texttt{lumen\_warehouse.sf\_origin} and \texttt{AH}.
  \item \textbf{Secrets:} generated server-side into
        \texttt{\textasciitilde/ahportal/app/.env.local}; staff accounts
        (\texttt{admin@}/\texttt{advisor@}/\texttt{manager@ttu.edu}) created
        with generated passwords.
  \item \textbf{Checklist pipe:} added the Lumen connector \emph{SF Checklist
        Items $\rightarrow$ AH} (id \texttt{3e0ca7b6}, incremental on
        \texttt{SystemModstamp}) so \texttt{TargetX\_SRMb\_\_Checklist\_Item\_\_c}
        lands in \texttt{AH.checklist\_items} --- 1.46M rows, all-applications.
  \item \textbf{Build \& run:} \texttt{npm ci \&\& npm run build}
        (webpack), PM2 \texttt{april-web} on :3210 plus the
        \texttt{ahportal-db-tunnel} SSH forward to 0003; \texttt{pm2 save}.
  \item \textbf{Mount:} nginx \texttt{location /sim/AHPortal/} in
        \texttt{chatbot.conf} $\rightarrow$ \texttt{127.0.0.1:3210} with
        forwarded-host headers (0001 already owns the public :443 cert).
  \item \textbf{Verified live:} portal 200 on the real cert, HTTPS login
        $\rightarrow$ session $\rightarrow$ admin stats over clone data, signup
        POST passes CSRF through the proxy, assets load under the prefix, boot
        sync confirmed idempotent through the tunnel (0 changes on re-run).
  \item \textbf{Before real students:} set \texttt{TWILIO\_*} for live SMS MFA
        (dev code \texttt{424242} disappears automatically once set) and point
        \texttt{BOT\_ENDPOINT} at the knowledge-base bot.
\end{enumerate}

\section{Staff Accounts \& Where They Live}

Sign in at \texttt{https://tosmonline0001.ttu.edu/sim/AHPortal/login}.

\begin{center}
\begin{tabularx}{\textwidth}{lXl}
\toprule
\textbf{Account} & \textbf{Password} & \textbf{Role} \\
\midrule
\texttt{admin@ttu.edu} & \texttt{zOKsKDcKBRYBe3} & superadmin \\
\texttt{advisor@ttu.edu} & \texttt{CTFR3z1MViZGwL} & advisor \\
\texttt{manager@ttu.edu} & \texttt{MwugOcgLjuV0F} & manager \\
\bottomrule
\end{tabularx}
\end{center}

While \texttt{TWILIO\_*} is unset, MFA challenges accept the dev code
\texttt{424242}. The same credentials live server-side at
\texttt{\textasciitilde/ahportal/STAFF\_LOGINS.txt} on tosmonline0001 ---
change them at first login and rotate the file copy if the document is shared.

\section{What Happens If Something Breaks}

\begin{center}
\begin{tabularx}{\textwidth}{lX}
\toprule
\textbf{Scenario} & \textbf{Result} \\
\midrule
Database unreachable & Pages error cleanly; no data is fabricated; the pool
reconnects automatically. \\
Twilio outage & OTP sends fail; the user sees an error; sessions already active
are unaffected. \\
Upload disk fills & New uploads refuse with an error; nothing else degrades. \\
Rate-limit false positive & The legit user waits out the \texttt{Retry-After}
window (minutes), not a ban. \\
Forgotten admin password & Another admin can reset it; a superadmin can update
the user row directly. \\
Server reboot & Sessions survive (they live in the database); the app restarts
via systemd. \\
\bottomrule
\end{tabularx}
\end{center}

\chapter{Open Items \& Known Trade-offs}

\section{Still To Receive / Decide}

\begin{itemize}[leftmargin=1.4em]
  \item \textbf{General Business requirements} --- the one program whose SF
        checklist isn't finalized; its applicants see ``requirements coming
        soon'' until TTU publishes them.
  \item \textbf{Twilio production keys} for live SMS MFA.
  \item \textbf{The custom chatbot endpoint} (\texttt{BOT\_ENDPOINT}) ---
        wire-in point is ready: the portal posts the message and user context,
        expects a reply string.
  \item \textbf{Database role details} for production (which account owns
        \texttt{AH}, what to grant).
  \item \textbf{Lead definition} --- currently ``all prospective-student
        contacts''; if the business definition of a lead is narrower, the
        filter is a one-line change.
\end{itemize}

\section{Honest Limitations}

\begin{itemize}[leftmargin=1.4em]
  \item \textbf{Chat polls} every 2.5 seconds rather than pushing instantly ---
        invisible to users at this scale, but worth upgrading to websockets if
        the staff queue grows busy.
  \item \textbf{Rate limits live in server memory} --- correct for one machine;
        if the app ever runs on two machines simultaneously they would need a
        shared store (the database table is a drop-in).
  \item \textbf{JWT role staleness} --- a staff member whose role is reduced
        keeps their old role until their next login (max 7 days). Immediate
        revocation happens by disabling the account or deleting their session
        row.
  \item \textbf{Full-disk encryption} of the database volume is an
        OS-level server task, not something the application can do for itself.
\end{itemize}

\chapter*{Glossary}
\addcontentsline{toc}{chapter}{Glossary}

\begin{description}[leftmargin=1.4em,style=nextline]
  \item[RFI] Request For Information --- the first contact a prospect makes
  (event signup, call, web form), which becomes a Salesforce lead.
  \item[RaiderConnect] TTU's official applicant portal
  (\texttt{texastech.my.site.com/online}) --- where the application itself is
  completed and where advisors push student documents into Salesforce.
  \item[Parity] The moment the warehouse shows Salesforce reflecting something
  the portal recorded (e.g.\ a pushed document whose checklist item is now
  received) --- only then does the portal treat it as final.
  \item[Lead] A prospective student in Salesforce who has not yet applied.
  \item[Salesforce mirror] A copy of our Salesforce data synchronized into our
  own database (\texttt{sf\_origin}) by the warehouse ETL --- read-only to us.
  \item[Clone (\texttt{AH.sf\_*})] The portal's own copy of just the
  TTU~Online slice of the mirror, refreshed every 5--30 minutes, with
  append-only history tables keeping every version of every row.
  \item[Sync engine] The in-app job that refreshes the clone on a timer or on
  demand from \texttt{/admin/sync}; advisory-locked so runs never overlap.
  \item[Confidence score] The 0--100 match strength from comparing email,
  phone, last name (including former names), and first name (including
  nicknames).
  \item[Projection] The generated columns stitched onto a linked SF contact:
  portal user, application status, status timeline, checklist progress.
  \item[Schema] A named compartment inside a database that keeps tables
  organized --- \texttt{sf\_origin} holds Salesforce data, \texttt{AH} holds
  ours.
  \item[Identity link] The recorded connection between a portal account and a
  Salesforce contact.
  \item[Outbox] The queue of pending changes a superadmin applies to Salesforce
  manually --- the bridge between the app and the system of record.
  \item[Hash] A one-way fingerprint of a password --- can verify a guess, can
  never reveal the original.
  \item[Pepper] A secret mixed into every password hash, stored only in server
  configuration.
  \item[MFA/OTP] Multi-Factor Authentication / One-Time Passcode --- the texted
  code that proves you also control your phone.
  \item[Session] The signed cookie that keeps you logged in between page loads.
  \item[Rate limit] A speed limit on requests that stops brute-force attacks.
  \item[CSRF] An attack where another website tricks your browser into
  submitting actions on our site --- blocked via origin checks.
  \item[CSP] Content-Security-Policy --- a whitelist telling the browser exactly
  which scripts and resources may run.
  \item[HSTS] The instruction that tells browsers ``HTTPS only, forever.''
  \item[AES-256-GCM] The encryption algorithm protecting uploaded documents at
  rest --- the same standard used for classified traffic.
  \item[Magic bytes] The identifying bytes at the start of a real file (a true
  PDF always begins \texttt{\%PDF}) --- used to detect faked uploads.
  \item[Least privilege] Every account gets only the access its job requires.
  \item[RBAC] Role-Based Access Control --- permissions follow the role, not the
  person.
\end{description}

\end{document}
